% context: markscheme (official solutions) for 2-6CW_Kinetics-review
% topic: kinematics review — distance vs displacement, average speed vs velocity,
%        average speed with a conversion to km/h, braking with a given negative a,
%        a dropped stone, a ball thrown straight up
% convention: answers in bold; "=" joins exact equalities, "\approx" introduces a
%        rounded value; each solution lists the knowns, then names the equation used.
%        Q2 follows the two-step rate rule: Step 1 the rate, Step 2 the conversion
\documentclass[12pt, twoside]{article}
%\documentclass[12pt, twoside]{article}
\usepackage[letterpaper, margin=1in, headsep=0.2in]{geometry}
\setlength{\headheight}{0.6in}
%\usepackage[english]{babel}
\usepackage[utf8]{inputenc}
\usepackage{microtype}
\usepackage{amsmath}
\usepackage{amssymb}
%\usepackage{amsfonts}
\usepackage[nomessages]{fp} %\FPeval{\var-name}{2*sin(pi/6)}
\usepackage{siunitx} %units in math. eg 20\milli\meter
\usepackage{yhmath} % for arcs, overparenth command
\usepackage{tikz} %graphics
\usetikzlibrary{quotes, angles, arrows, arrows.meta}
\usepackage{graphicx} %consider setting \graphicspath{{images/}}
\usepackage{parskip} %no paragraph indent
\usepackage{enumitem}
\usepackage{multicol}
\usepackage{venndiagram}

\usepackage{fancyhdr}
\pagestyle{fancy}
\fancyhf{}
\renewcommand{\headrulewidth}{0pt} % disable the underline of the header
\raggedbottom
\hfuzz=2mm %suppresses overfull box warnings

\usepackage{hyperref}
\usepackage{float}
\setlength{\parskip}{0.2\baselineskip}

\title{Physics}
\author{Chris Huson}
\date{November 2026}

\fancyhead[LE]{\thepage}
\fancyhead[RO]{\thepage}
\fancyhead[LO]{La Scuola d'Italia / Huson / Physics \\* 5 November 2026}

\begin{document}

\subsubsection*{2.6 Classwork: Kinematics Review Solutions}

\begin{enumerate}[itemsep=0.5cm]

\item Hallway walk, east positive, origin at the door.

  \begin{center}
  \begin{tikzpicture}[scale=0.6]
    \draw[<->, thick] (-6.6,0) -- (8.6,0) node[right]{$x$ (m)};
    \foreach \x in {-6,...,8} \draw (\x,0.12) -- (\x,-0.12);
    \foreach \x in {-6,-4,...,8} \node[below] at (\x,-0.12) {\small \x};
    \node[above] at (0,0.1) {\scriptsize door};
    \draw[->, very thick] (0,0.7) -- (6,0.7) node[midway, above]{\small $+6.0$ m};
    \draw[->, very thick] (6,1.6) -- (-4,1.6) node[midway, above]{\small $-10.0$ m};
  \end{tikzpicture}
  \end{center}

  \begin{enumerate}[itemsep=0.3cm]
    \item $x = 0 + 6.0 - 10.0 = \mathbf{-4.0~\textbf{m}}$, that is, 4.0 m west of the
      door.
    \item Distance = $6.0 + 10.0 = \mathbf{16.0~\textbf{m}}$. \quad Displacement =
      $-4.0 - 0 = \mathbf{-4.0~\textbf{m}}$, or \textbf{4.0 m west}.

      \emph{Distance counts every meter walked and has no direction. Displacement is
      only where the student ended up, measured from the start.}
    \item Total time = $3.0 + 5.0 = 8.0$ s.

      average speed $= \dfrac{16.0~\text{m}}{8.0~\text{s}} = \mathbf{2.0~\textbf{m/s}}$
      \qquad
      $\bar{v} = \dfrac{-4.0~\text{m}}{8.0~\text{s}}
      = \mathbf{-0.50~\textbf{m/s}}$, or \textbf{0.50 m/s west}

      \emph{A common slip is to average the two legs' velocities,
      $\frac{2.0 + (-2.0)}{2} = 0$; that ignores that the legs take different
      times. Divide the total displacement by the total time.}
  \end{enumerate}

\item Usain Bolt, 100 m in 9.58 s. The race distance is exact, so the answer takes
  three significant figures from 9.58 s.
  \begin{enumerate}[itemsep=0.3cm]
    \item Step 1, the rate:
      average speed $= \dfrac{100~\text{m}}{9.58~\text{s}} = 10.438\ldots
      \approx \mathbf{10.4~\textbf{m/s}}$

      Step 2, convert, carrying the unrounded value:

      $10.438\ldots~\dfrac{\text{m}}{\text{s}} \times
      \dfrac{3600~\text{s}}{1~\text{h}} \times \dfrac{1~\text{km}}{1000~\text{m}}
      = 37.578\ldots \approx \mathbf{37.6~\textbf{km/h}}$

      \emph{Also accept $\times 3.6$ in Step 2. Rounding first, $10.4 \times 3.6 =
      37.44 \approx 37.4$ km/h, loses the third figure --- the reason to carry
      $10.438\ldots$}
    \item 12.3 m/s is an \textbf{instantaneous} speed. \textbf{He starts from rest},
      and the first few seconds are spent speeding up, so the average over the whole
      race is lower than the top speed.

      \emph{Credit any answer that names the start from rest (the acceleration
      phase). Also accept that he eases off slightly at the end.}
  \end{enumerate}

\item Braking car. Knowns: $v_0 = 17.9$ m/s, $v = 0$, $a = -6.4$ m/s$^2$. The
  acceleration has two significant figures, so the answers do too.
  \begin{enumerate}[itemsep=0.3cm]
    \item Want $t$. Use $v = v_0 + at$, so $t = \dfrac{v - v_0}{a}$:

      $t = \dfrac{0 - 17.9~\text{m/s}}{-6.4~\text{m/s}^2} = 2.796\ldots
      \approx \mathbf{2.8~\textbf{s}}$
    \item Want $x - x_0$; use $v^2 = v_0^2 + 2a(x - x_0)$:

      $0 = (17.9~\text{m/s})^2 + 2(-6.4~\text{m/s}^2)(x - x_0)$, \quad so

      $x - x_0 = \dfrac{-320.41~\text{m}^2/\text{s}^2}{-12.8~\text{m/s}^2}
      = 25.03\ldots \approx \mathbf{25~\textbf{m}}$

      \textbf{Yes, it agrees}: 25 m and 24 m differ by about 1 m, which is within
      the precision of a two-figure acceleration.

      \emph{Also accept the average velocity route: $\bar{v} = \frac{17.9 + 0}{2}
      = 8.95$ m/s, and $8.95 \times 2.796\ldots \approx 25$ m. A negative distance
      means the minus sign on $a$ was dropped in one place but not the other.}
  \end{enumerate}

\item Dropped stone. Knowns: $v_0 = 0$, $t = 2.00$ s, $a = +9.80$ m/s$^2$ (down
  positive).
  \begin{enumerate}[itemsep=0.3cm]
    \item Want $v$. Use $v = v_0 + at$:

      $v = 0 + (9.80~\text{m/s}^2)(2.00~\text{s}) = \mathbf{19.6~\textbf{m/s}}$
      downward
    \item Want $y - y_0$. Use $y = y_0 + v_0 t + \frac{1}{2}at^2$ with $v_0 = 0$:

      $y - y_0 = \tfrac{1}{2}(9.80~\text{m/s}^2)(2.00~\text{s})^2
      = 4.90 \times 4.00 = \mathbf{19.6~\textbf{m}}$

      \emph{The two answers share the digits 19.6 only because $t = 2.00$ s; one is a
      speed and one a height. Also accept $\bar{v}\,t = 9.80~\text{m/s} \times
      2.00~\text{s}$. The usual slip is forgetting the $\frac{1}{2}$ (39.2 m) or not
      squaring the time (9.80 m).}
  \end{enumerate}

\item Ball thrown up. Knowns: $v_0 = +9.80$ m/s, $a = -9.80$ m/s$^2$, $y_0 = 0$ at
  the hand (up positive).
  \begin{enumerate}[itemsep=0.3cm]
    \item \textbf{Velocity $\mathbf{v = 0}$} for an instant. \textbf{Acceleration
      $\mathbf{a = -9.80}$ m/s$^\mathbf{2}$}, downward, the same as everywhere else
      in the flight.

      \emph{The common wrong answer is $a = 0$ at the top. If the acceleration were
      zero there, the velocity would stay zero and the ball would hang in the air.}
    \item Want $t$ when $v = 0$. Use $v = v_0 + at$, so $t = \dfrac{v - v_0}{a}$:

      $t = \dfrac{0 - 9.80~\text{m/s}}{-9.80~\text{m/s}^2} = \mathbf{1.00~\textbf{s}}$

      \emph{The ball loses 9.80 m/s every second and starts with 9.80 m/s.}
    \item Want $y - y_0$. Use $\bar{v} = \dfrac{v + v_0}{2} = \dfrac{0 + 9.80}{2}
      = 4.90$ m/s:

      $y - y_0 = \bar{v}\,t = 4.90~\text{m/s} \times 1.00~\text{s}
      = \mathbf{4.90~\textbf{m}}$

      \emph{Also accept $v^2 = v_0^2 + 2a(y - y_0)$:
      $y - y_0 = \frac{0 - (9.80)^2}{2(-9.80)} = \frac{-96.04}{-19.6} = 4.90$ m.}
    \item \textbf{2.00 s.} The fall back down takes as long as the rise:
      $2 \times 1.00~\text{s} = 2.00$ s.

      \emph{It comes back at 9.80 m/s downward, $v = 9.80 + (-9.80)(2.00) = -9.80$
      m/s. Also accept setting $y = 0$ in $y = v_0 t + \frac{1}{2}at^2$:
      $0 = (9.80 - 4.90\,t)\,t$, so $t = 2.00$ s.}
  \end{enumerate}

\end{enumerate}
\end{document}
